\chapter{Theoretical Background}\label{ch:theory}
\section{Defining ``Cult''}\label{sec:defining_cult}
\subsection{Current definitions}\label{subsec:current_definitions}
Current dictionaries seems to give the same main definition - and the one that this research will focus on: \enquote{\textbf{\cultbelief{}} \textit{noun}: religious group, often living together, whose beliefs are considered extreme or strange by many people. 
\textit{Their son ran away from home and joined a religious cult.}} in \autocite{Cult2026CambridgeDefinition} or \enquote{a small group of people who have extreme religious beliefs and who are not part of any established religion} in \autocite{CultNounDefinitionOxford}

Usually this definition is also complemented by two other possible definitions, which I will henceforth refer to as \cultworship{} and \cultfame{} respectively:

\textbf{\cultworship{}} as a set of belief, still linked to religion or spirituality, for example: \enquote{\textbf{cult} \textit{noun}:\ a particular system of religious belief. \textit{The Roman cult of Mithras admitted only men.}}

\textbf{\cultfame{}} something more akin to fame, popularity, \enquote{\textbf{cult} \textit{noun}:\ someone or something that has become very popular with a particular group of people. \textit{the cult of celebrity}}, and in their american english dictionary: \enquote{a system of religious belief, esp. one not recognized as an established religion, or the people who worship according to such a system of belief.} This definition is also usually complemented by an adjective \enquote{\textbf{cult} \textit{adjective}:\ liked very much by a particular group of people. \textit{the cult of celebrity}}

\subsection{Etymology}
These different definitions match the alleged etymology. 
\cultworship comes first, in the 1610 with the obsolete sense of ``worship, hommage'', then in the sense defined above in 1670s from French \textit{culte}, itself derive from Latin \textit{cultus} originally closer to a ``care, labor'' sense.  The word is rarely used until the 1900s where it starts to get an understanding closer to refer to \cultbelief: \enquote{ancient or primitive systems of religious belief and worship}\autocite{CultEtymonline}

\subsection{A contested concept}
\subsubsection{Cults in the modern era}
The 1900s definition of the word `cult' opposes contemporary religious practices to older, non-legitimate ones. This aspect of legitimacy can easily be abstracted away from its temporal aspect, into a general negative connotation. Put more simply, Rawson defines: \enquote{\textbf{cult.} An organized group of people, religious or not, with whom you disagree.} \autocite[pp.~106]{rawsonWickedWordsTreasury1989}. Barker would go as far as saying that \enquote{To label a movement a cult can be to suggest that it is a dangerous pseudo-religion with satanic overtones which is likely to be involved in financial rackets and political intrigue, to indulge in unnatural sexual practices, to abuse its women and children, and to use irresistible and irreversible brainwashing techniques in order to exploit its recruits. Furthermore, it is implied, it may well resort to violence, perform various criminal activities and, possibly, persuade its members to commit mass suicide.}, before insisting on a strict separation between \enquote{definition} and \enquote{discovery}, moving away from considering this movements with assumptions to better be able to study them.

This negative-connotation axis is only one line along which the concept is contested; the rest of this section traces two others -- the proliferation of alternative terminology attempting to escape it (\enquote{Alternative Terminology}, below), and the question of whether the phenomenon should be bounded by religion at all (\enquote{Beliefs beyond religion}, below) -- before closing by noting that the sociological definitional problem is itself productive of new terms, rather than resolved by any one of them.

\subsection{Alternative Terminology}
Other words have been used to describe the same phenomena: \enquote{The movements or groups with which this paper is concerned have been labelled `cults', `sects'\footnote{
	Often, cults and sects seems to be used interchangeably, although their distinction has been formalised in 1979: sects are schismatic breakaways from an existing religious body (they have a lineage), while cults are novel, typically formed around a charismatic leader with no prior institutional tie. \autocite[pp.125]{starkChurchesSectsCults1979}.
}
, `new religious movements (NRMs)', `minority religions', `alternative religions', `spiritual or faith communities'}.\autocite[pp.~236]{barkerNotsonewReligiousMovements2014}. 

\subsubsection{New Religious Movements (NRMs)}
The concept of \textbf{New Religious Movements(NRMs)} –- also rarely but at times called `alternative religions' -- is now also widely used by reseachers since the 70's,because it offers a non-judgemental or pejorative connotations, as opposed to `sects' and `cults'. It is usually preferred to terms like `deviant religions', `extremist religious movements', `fanaticsm'.
NRMs are also conceptually contested for multiple reasons. First they restrict these movements to a religious aspect which is by itself already difficult to encircle.\footnote{\enquote{If we confine it to belief in a God or gods, then much of Buddhism cannot be called a religion. Several groups that are termed NRMs insist they do not want to be called religious}\autocite[pp.237]{barkerNotsonewReligiousMovements2014}}. Second, movement can play with the institutional weight of the \textit{religion} category for their own interest, whether seeking protection by law for religion (Realians), getting tax breaks(Scientology), or fleeing the label to be able to be taught in schools(Transcendental Meditation).
It is admitted that \enquote{scholars have tended to accept a relatively comprehensive understanding of religion, following the suggestion by Tillich (1957, 1–4) that faith, including religion, is that which offers answers to questions of ultimate concern, such as \textit{Is there a God?}, \textit{What is the purpose of life?}, \textit{What happens at death?}}\autocite[pp.237]{barkerNotsonewReligiousMovements2014}

The term movement can also be contested, as well as their newness: some of these movement claim that they have been around for centuries, especially when they branched out of an older established religion. A key aspect would be provided by \citeauthor{barkerNotsonewReligiousMovements2014}, defining NRMs as \enquote{first-generation religions}\autocite[pp.239]{barkerNotsonewReligiousMovements2014}.

\subsubsection{Beliefs beyond religion}

Since religion is also an ambiguous concept, some scholars even talk about a wider category they would call ``extreme beliefs''. 
Extreme can be understood in different ways. First, when looking at content of the beliefs, as a deviation from the norm\footnote{Extreme beliefs as a deviation from the norm are obviously critisizable, mainly because of the contextual nature of norms. For example a legal policy can be democratically voted somewhere and be rejected somewhere else.}. 
More subtly, J.M. Berger, frames extremism as \enquote{the belief that an in-group can never be healthy or successful unless it is engaged in hostile action against an out-group.}\autocite{bergerExtremism2018}
Another aspect of extreme beliefs could be their `epistemic resiliance' with rigid convictions that resist counter-evidence and often form the core of movements like extremism, fundamentalism, and radical conspiracy theorizing.\autocite{bortolottiExtremeBeliefsUnshakable2025}

Instead of widening the category, some proposed to use religion as a framework to understand the variety of movement with similar behaviours,this is the quasi-religious hypothesis: conspiracy theories share content, form, and function with institutionalized religious belief -- meaning-making, agency-detection, moral polarization, in-group/out-group construction -- without being identical to or causally linked with religiosity, and have been framed as a possible ``surrogate for God'' in secularizing societies \autocite[pp.2]{franksConspiracyTheoriesQuasireligious2013}. 

Generally speaking, the problem of a sociological definition has created multiple terms, and \enquote{Different types of scholars can be concerned about different aspects of the phenomena they are studying}.\autocite[pp.239]{barkerNotsonewReligiousMovements2014}.


\section{Family resemblance and prototype theory}\label{sec:prototype_theory}
Wittgenstein, using the example of the concept of \textit{game}, arrived at the concept of \textbf{family resemblance}: a word can pick out a class of objects that share no single defining feature, only a network of overlapping similarities -- the way members of a family resemble one another without any one trait being common to all of them.\autocite{wittgensteinPhilosophischeUntersuchungenPhilosophical2010}, which goes again a conception where a term would point to a single defining essence.

In cognitive psychology, this idea was subsequently developed into \textbf{prototype theory}.Rosch and Mervis proposed that categories can possess an internal structure: their members differ in degree of typicality, and the most prototypical members tend to share many attributes with other members of their own category while sharing fewer attributes with members of contrasting categories \autocite{roschMervisFamilyResemblances1975}.

Categorisation, on this account, is not necessarily based on checking each instance against a fixed set of necessary and sufficient conditions. Instead, judgments of category membership may be influenced by graded similarity to salient or central examples, producing more typical and less typical members as well as borderline cases.


\section{Relevant Epistemologies}
The sociological attempt to find a conceptual frame for these movements, shows that abstracting away from a judgement of the beliefs a group holds is also, necessarily, a shift of focus toward social structure: away from \textit{what} a group believes and toward \textit{how} it is organised, and how it relates to the society around it.

In order to understand the social structure of these movements, this papers selected three different main epistemologies around the word \enquote{cult} and its related terms.

\subsection{Interviews}
Directly building on the prototype theory outlined above (\ref{sec:prototype_theory}), the interviews are this thesis's way of accessing lay usage of \enquote{cult}: rather than asking participants to state a definition, the interview protocol (\ref{sec:interview_protocol}) asks what comes to mind first when the word is heard, then asks the participant to justify that specific example. The elicited object is a prototype -- a concrete exemplar treated as central to the category -- rather than an attempted definition, on the assumption that this is closer to how the category is actually organised cognitively for most speakers, and easier to elicit reliably in a short interview than a definition would be. This exemplar is what the Methods chapter's interview prototype layer (\ref{sec:vectorising_prototypes}) later places in the shared space. This framework was then used significantly in lexical semantics amd cognitive lingustics. 

\subsection{Scholarly Epistemology}
- Systematic academic theorising about cults/NRMs, accumulated and revised over decades of case-study research (sociology of religion, NRM studies, psychology of extreme belief) -- the epistemology already surveyed in \ref{sec:defining_cult} above, rather than introduced fresh here.
- Explicitly bracketed off from any single agreed definition: the field's own move has been away from adjudicating \enquote{is X really a cult} and toward describing the criteria different scholars actually use (\enquote{definition}, not \enquote{discovery}, \ref{sec:defining_cult}).
- Internally heterogeneous rather than a single school: competing and partially-overlapping frameworks (the superseded brainwashing paradigm, NRM studies, the quasi-religious hypothesis, the wider extreme-beliefs framing) coexist rather than converge.
- The largest and most heterogeneous of the three epistemologies by volume -- decades of publication, unlike MIVILUDES's narrower operational output or lay usage, which exists nowhere in documented form until this thesis's own interviews produce it.
- Operationalised via a reproducible corpus-construction protocol (keyword search over OpenAlex, then screening and disciplinary stratification) rather than a single canonical reading list, precisely because no such canonical list exists -- see the corpus construction protocol (Methods, \ref{sec:corpus_construction_protocol}) and how it is vectorised (\ref{sec:vectorising_scholarly_work}).

\subsection{Legal Epistemology}
% Quick history, definition of "sectarian drift" (dérive sectaire) and its
% indicator lists.
\subsubsection{Why MIVILUDES Specifically}
- France is close to a limit case for this comparison: its la\"icit\'e model is one of the few legal systems that formally distinguishes \enquote{true religions} from \textit{sectes} and gives the state a mechanism to act on that distinction, whereas most other systems -- the US Establishment Clause is the clearest example -- have no legal category for \enquote{cult} at all, and treat all religions as formally equal (see the Introduction, \enquote{Legitimacy is legally/culturally constructed}). A legal/administrative epistemology of \enquote{cult} is easiest to study where one has actually been built into state machinery.
- MIVILUDES itself already moved away from labelling groups \enquote{cults} and toward \textit{dérive sectaire} (sectarian drift): a consequence-based, structural criterion rather than a content-based one (below, \enquote{Sectarian Drifts}) -- which is exactly the kind of structural, non-belief-based operationalisation this thesis is oriented toward across all three epistemologies.
- It gives this thesis a single, bounded, well-documented source rather than a comparative-law synthesis across many jurisdictions: a publicly available activity report, a public list of seventeen official criteria, and a defined intake process (below), all usable directly rather than needing to be reconstructed from case law.
- Its own reports already quantify and categorise what they receive (the \enquote{Dominant themes} breakdown below), and its seventeen criteria are official, fixed, and already the right shape to embed and compare geometrically (Methods, \ref{sec:vectorising_miviludes}) -- unlike the scholarly corpus, which offers no single agreed set of dimensions to anchor to in advance.

\subsubsection{History of the MIVILUDES}
- Before MIVILUDES: the mission's predecessor, MILS (Mission interministérielle de lutte contre les sectes, est. 1998), drew international criticism for religious-freedom overreach; it was abolished and replaced by MIVILUDES under decree n°2002-1392 (28 Nov 2002) \autocite{decret2002-1392}.
- The moral panic (and what some scholars call ``anti-cult information terrorism'') around new religious movements spiked sharply after the murder-suicide of 53 members of l'Ordre du Temple Solaire in October 1994 (Switzerland and Quebec) \autocite[pp.4]{adeliyantousUsingLawLimit2023}.
- Most EU countries judged existing criminal sanctions sufficient to punish dangerous NRM conduct; France instead built a distinct approach, managing the issue through governmental measures coordinated with non-governmental anti-cult associations \autocite[pp.4]{adeliyantousUsingLawLimit2023} -- this is what eventually produced MIVILUDES.
- Stanley Cohen's concept of \textit{moral panic} is the companion idea here: a ``widespread fear, often irrational, of a person or phenomenon perceived to threaten social values'' \autocite{cohenFolkDevilsMoral2011}.

\subsubsection{Sectarian Drifts}
The creation of MIVILUDES and the creation of the concept of Sectarian Drifts implies to go away from pointing the finger at groups and identify them as cult, and rather focus on the negative consequences they have socially.
Indeed, out of all the reports (\textit{signalements} in french) made to the MIVILUDES in the period 2022-2024, the following `Dominant themes' were reported:

\begin{itemize}
	\item 37\% Health and well-being.
	\item 35\% Worship and spirituality.
	\item 13\% Training, jobs, and finances
	\item 6\% Plot theory, separatism, and other forms of extreme beliefs (like survivalisme)
	\item 9\% Other
\end{itemize}\autocite[pp.8]{missioninterministerielledevigilanceetdeluttecontrelesderivessectairesmiviludesRapportDactiviteMiviludes2024}


\subsubsection{How MIVILUDES Operates}
\enquote{Miviludes receives reports and requests for information from individuals and legal entities—both private and public—who contact the agency to seek its expertise.}\autocite[pp.10]{missioninterministerielledevigilanceetdeluttecontrelesderivessectairesmiviludesRapportDactiviteMiviludes2024}
After receving them the report is treated by a specialised counselor who redirects the person to competent institutions - or to the prosecutaor if there seems to be a legal infraction - and  other entities, like victim organisations, for example. 
The MIVILUDES cannot conduct their own investifations, and have no legal power.


The MIVILUDES also publishes reports that give insight into its own work \autocite{missioninterministerielledevigilanceetdeluttecontrelesderivessectairesmiviludesRapportDactiviteMiviludes2024}.



\section{Semantic Spaces, Conceptual Spaces, latent spaces and embeddings}
% Background on semantic/conceptual spaces as a representational framework.


\subsection{Semantic Spaces}
A semantic space represents meaning geometrically: a word, phrase, or document is placed as a point (or a region) in a space whose dimensions are derived from patterns of use, so that geometric proximity between two points stands in for semantic relatedness between what they represent. The underlying assumption is the distributional hypothesis: a word's meaning is reflected in the contexts it occurs in, so that words appearing in similar contexts tend to be similar in meaning \autocite{harrisDistributionalStructure1954, firth1957synopsis}. \enquote{Semantic space} is the general term for this family of representations; conceptual spaces, latent spaces, and embeddings, discussed in turn below, are three more specific ways of constructing and interpreting such a space, distinguished by how their dimensions are obtained and by what can legitimately be inferred from a point's position within them.

\subsection{Conceptual Spaces}
Gärdenfors's conceptual spaces \autocite{gardenforsConceptualSpacesGeometry2000, gardenforsConceptualSpacesGeometry2004} are a specific, more constrained proposal within this broader tradition: concepts are represented as convex regions within a space built from a small number of interpretable, often perceptual or cognitive, quality dimensions (e.g.\ hue, or a temperature scale for the concept \textit{hot}), typically with a prototype-like centre (\ref{sec:prototype_theory}) rather than an arbitrary boundary.

A caution is worth stating explicitly here, since the two are easy to conflate: in Gärdenfors's strict sense, concepts are regions in a space with interpretable domains and, often, prototype-like centres; in an ordinary embedding (below), a word or passage is instead a single point or vector whose individual dimensions normally lack any stable, direct interpretation of their own. For this reason, this thesis calls the spaces it builds \enquote{distributional semantic spaces} or \enquote{learned representational spaces} throughout, and reserves \enquote{conceptual space} specifically for this stricter, Gärdenfors sense -- used here only as a point of theoretical comparison, not as a claim that the spaces built later in this thesis are conceptual spaces in his sense. Whether they could be related more rigorously to conceptual spaces proper -- for instance by testing the convexity of a candidate region, or by attempting to recover interpretable dimensions from it -- is left as a direction for future work, taken up again in the concluding chapter.



\subsection{Latent Spaces}
A latent space is a learned representational space whose variables are not directly observed, but inferred because they help explain, predict, compress, or generate the data a model was trained on; it is defined by a model architecture, a training objective, and a body of training data, rather than by hand-specified dimensions. It differs from a semantic space in general (above) in status rather than in kind: \enquote{latent} does not guarantee semantic interpretability, and does not hold some kind of truth about the reality of the concepts it manages -- it designates an emergent structure that is useful for explaining or predicting observations, a compressed statistical structure recovered from language itself rather than imposed on it in advance.


“Latent” names the status and mode of learning of a representation, not its semantic content. A latent variable is unobserved: it is inferred because it helps explain, predict, compress, or generate observed data. Thus, an embedding can be latent, but not all embeddings are necessarily latent in the same way—for example, an explicit PPMI word–context row is a vector representation with interpretable context coordinates, even if it is later compressed through SVD.

\subsection{Embeddings}
An embedding is a learned mapping from a discrete unit of language -- a word, a phrase, a passage -- to a dense numerical vector, typically produced by a neural model trained on a large text corpus; embeddings are the concrete, computational instance of a semantic space (above) that this thesis actually builds and measures. They are not direct measurements of \enquote{meaning}, but empirically and mathematically motivated relational representations: their geometry can be treated as evidence about patterns of meaning or use when the unit of inference is the relative configuration -- similarities, neighborhoods, cluster separations, and validated directional contrasts -- rather than the absolute value of a coordinate, or a naïve reading of a single vector dimension. This use of embedding geometry as evidence, and its methodological caveats, is by now well documented in the wider NLP literature \autocite{worthWordEmbeddingsSemantic2023, simhiInterpretingEmbeddingSpaces2023}.

Embedding geometry is an empirically testable operationalisation of semantic relatedness. Because the training objective couples vectors to contextual prediction or association, and because the resulting relational configurations reproduce independently-assessed semantic, syntactic, and social regularities, distances, neighborhoods, and certain validated directions can be used as evidence about structured patterns of meaning and use.

The strongest argument is cumulative rather than metaphysical. Geometry is meaningful because models are trained to preserve or predict contextual relations; because the learned spaces reproduce independently assessed semantic relations; because some geometric regularities can be derived from corpus statistics; and because the same geometry supports prediction, retrieval, category separation, and controlled transformations. That is evidence for construct validity, not proof that vectors reveal an unmediated ontology of concepts.



\section{Conclusion}

This chapter has argued three things. First, that \enquote{cult} is not a term with a single, agreed content: dictionary, etymological, and scholarly treatments converge on a rough sense -- a religious or quasi-religious group whose beliefs or practices are considered extreme -- but that sense turns out to rest on a legitimacy judgement rather than a content-based one (\ref{sec:defining_cult}), and immediately fragments into competing terminology (NRM, sect, extreme belief, quasi-religion) each attempting to manage that same underlying problem differently: how to describe such groups without simply importing a verdict on them.

Second, that this problem is best approached not by adjudicating between definitions, but by comparing how the term is actually used across three epistemologies that relate to it differently -- scholarly theorising, MIVILUDES's operational and legally-consequential framework, and lay, prototype-based usage -- treating none of the three as authoritative over the others, and treating the differences in how much each has been documented as a fact about the epistemologies themselves rather than a flaw in the comparison.

Third, that comparing these epistemologies at all requires putting them into a shared representational space, and that embeddings -- an empirically grounded, if epistemically modest, operationalisation of semantic relatedness -- offer exactly that: a way to place expressions of \enquote{cult} drawn from all three epistemologies into one geometry, where relative position, not any single coordinate, is the unit of evidence. The Methods chapter builds this space; the Results chapter puts it to work on the questions posed in the Introduction.

